\documentclass[../DoAn.tex]{subfiles}
\begin{document}

\newcommand{\anh}[3]{%
  \IfFileExists{Hinhve/#1}%
    {\includegraphics[width=#2]{#1}}%
    {\fbox{\parbox[c][#3][c]{#2}{\centering\footnotesize Ảnh chụp đặt ở đây:\\\texttt{\detokenize{#1}}}}}}

Chương này trình bày cách hiện thực các thiết kế ở Chương~\ref{chapter:Design}, theo thứ
tự môi trường, phần cứng, mô hình, firmware, backend, giao diện web rồi đường đưa hệ thống
lên máy chủ. Bản đang chạy khác thiết kế ở một điểm lớn: máy lắp trên breadboard, chưa có
mạch đế in, chưa có vỏ và thanh chắn. Mục~\ref{section:4.2} nói rõ phần này.

\section{Môi trường triển khai}
\label{section:4.1}

Mã nguồn nằm trong một kho chung, và mọi thao tác đi qua một Makefile ở gốc kho: cài đặt,
sinh mã từ hợp đồng, huấn luyện, lượng tử hóa, build, nạp và phát hành. Một số đo trong sổ
đo vì thế đi kèm đúng một dòng lệnh để lặp lại. Bảng~\ref{tab:moi_truong} liệt kê công cụ
và phiên bản. Máy phát triển là một laptop chạy WSL2, vừa build vừa huấn luyện. WSL chỉ có
10~GB RAM, nên các bước nặng của thang lượng tử hóa chạy trong một scope systemd có trần bộ
nhớ: bước nào vượt trần thì chỉ bước đó chết, máy không sập theo.

\begin{table}[!htb]
\centering
\small
\caption{Công cụ và phiên bản.}
\label{tab:moi_truong}
\begin{tabular}{|L{2.8cm}|L{10.8cm}|}
\hline
\textbf{Khối} & \textbf{Công cụ} \\
\hline
Máy phát triển & Laptop chạy WSL2, RAM 10~GB, GPU RTX 3050 Laptop \\
\hline
Mạch đế & KiCad 10.0.6, script Python sinh sơ đồ và mạch in \\
\hline
Mô hình & Python 3.12, PyTorch 2.13 với CUDA 13.0, ONNX, ESP-PPQ 1.3.11 \\
\hline
Firmware & ESP-IDF 6.0.2, esp-dl 3.3.11, esp-tflite-micro 1.4.0, esp-nn 1.3.2, esp32-camera 2.1.7, LittleFS 1.22.3 \\
\hline
Backend & Node.js 24, NestJS 12, Prisma 7, BullMQ 6, Socket.IO 4, TypeScript 6 \\
\hline
Frontend & Next.js 16, React 19, Kumo 2.14.0, TailwindCSS 4, TanStack Query 5, next-intl 4 \\
\hline
Android & Trusted Web Activity, androidbrowserhelper 2.6.2, SDK đích 36, SDK thấp nhất 21 \\
\hline
Máy chủ & VPS 1 vCPU, 2~GB RAM, Ubuntu 24.04, Docker Compose, Traefik v3, EMQX 6.3.1, PostgreSQL 16, Redis 7 \\
\hline
Phát hành & GitHub Actions, kho ảnh container của GitHub, Vercel \\
\hline
\end{tabular}
\end{table}

Phiên bản được ghim chặt ở những chỗ có hợp đồng hai phía. ESP-PPQ 1.3.11 và esp-dl 3.3.11
đi thành cặp, vì mỗi bản vá đồ thị ở Mục~\ref{section:4.3} phải khớp kernel trên chip từng
số nguyên. Kumo ghim đúng 2.14.0 vì thư viện này lên bản liên tục.

\section{Lắp ráp phần cứng}
\label{section:4.2}

Máy hiện lắp trên breadboard bằng dây cắm, theo bảng chân ở Phụ lục~\ref{appendix:chan},
như Hình~\ref{fig:ban_lap}. Mạch đế ở Mục~\ref{section:3.3.3} đã xong thiết kế, đạt ERC và
DRC, đã xuất tệp Gerber nhưng chưa đặt in, nên các luật đi dây của mục đó chưa được thử trên
đồng thật. Máy chưa có vỏ và chưa có thanh chắn: servo SG90 quay trần trên bàn, góc 0° là
đóng, 90° là mở. Những chỗ đồ án viết mở thanh chắn vì thế chỉ tới mức servo quay đúng góc.
Mọi số đo ở Chương~\ref{chapter:Experiment} lấy trên bản lắp này.

\begin{figure}[!htb]
\centering
\anh{anh_ban_lap.jpg}{0.8\textwidth}{6cm}
\caption{Máy chấm công lắp trên breadboard.}
\label{fig:ban_lap}
\end{figure}

Bản đang chạy lấy hai đường 5~V của Bảng~\ref{tab:nguon} từ một khối pin nằm ngoài mạch,
Hình~\ref{fig:khoi_pin}. Khối này gồm hai pin 18650 3,7~V 3.300~mAh nối tiếp qua mạch bảo
vệ 2S, một công tắc và hai module hạ áp MP1584EN, mỗi module nuôi một đường. Hộp pin đo được
8,3~V, hai module chỉnh ở 5,11--5,12~V. Khối không có mạch sạc nên pin tháo ra sạc ngoài.
Trước khi lắp vào máy, module của đường 1 được thử riêng bằng ứng dụng phát tiếng. Lấy điện
cho bộ khuếch đại qua chân của board thì loa ra tiếng như khi dùng USB, còn kéo một dây 5~V
riêng tới bộ khuếch đại trên breadboard thì loa rè to rồi mất tiếng. Đồ án nghi đường mass
của dây cắm nhưng chưa chứng minh được, và cũng chưa đo thời gian chạy của một lần sạc.

\begin{figure}[!htb]
\centering
\anh{anh_khoi_pin.jpg}{0.6\textwidth}{5cm}
\caption{Khối pin 2S và hai module hạ áp.}
\label{fig:khoi_pin}
\end{figure}

Ngoại vi được đưa lên lần lượt, mỗi con có một ứng dụng kiểm thử riêng chạy trên board trước
khi ghép vào firmware. Cách làm này lộ ra một loại lỗi mà build và kiểm thử trên máy tính
không bắt được: hàm báo thành công trong khi phần cứng chưa làm được việc.
Bảng~\ref{tab:loi_board} ghi các lỗi điển hình. Nghiệm thu màn hình vì thế phải có người
nhìn kính, vì bộ kiểm thử panel có 9 ca thì 8 ca chỉ kiểm mã trả về.

\begin{table}[!htb]
\centering
\small
\caption{Các lỗi chỉ lộ ra trên board.}
\label{tab:loi_board}
\begin{tabular}{|L{4.3cm}|L{5.0cm}|L{4.2cm}|}
\hline
\textbf{Hiện tượng} & \textbf{Nguyên nhân} & \textbf{Cách xử lý} \\
\hline
Lần ghi I\textsuperscript{2}C đầu bị NACK, chỉ sau khi rút điện cắm lại & Hàm khởi tạo bus trả về ở mốc 4~ms, PCF8574 trả lời ở 5~ms, VL53L1X ở 17~ms. Reset chip không tái hiện được vì ngoại vi còn điện & Thăm dò tới khi PCF8574 trả lời, trần 500~ms \\
\hline
Ngắt VL53L1X đủ 20/20 lần nhưng lần đọc nào cũng chưa có dữ liệu & Cảm biến mặc định tích cực mức cao, driver bắt cạnh xuống nên đọc đúng lúc dữ liệu vừa bị xóa & Đặt tích cực mức thấp ngay trong hàm khởi tạo \\
\hline
Màn hình tối hẳn sau khi thêm một phép kiểm cảm ứng & Phép kiểm trả lỗi, \texttt{ESP\_ERROR\_CHECK} reset chip khoảng mỗi 150~ms & Chia ngoại vi thành nhóm dừng máy và nhóm chạy tiếp (Mục~\ref{section:3.5.6}) \\
\hline
Màn sáng trơ, mọi lệnh SPI báo thành công, log vẫn 14~fps & Một chân dây SPI lỏng. Bus này một chiều nên firmware không có cách biết & Tô đỏ, lục, lam để mắt kiểm tra, rồi cắm lại dây \\
\hline
GT911 trả lời ở 0x14 khi chưa chạy trình tự reset & Chân reset nằm trên PCF8574 và tự lên cao lúc cấp nguồn, nên địa chỉ do cuộc đua quyết định & Luôn chạy trình tự reset, đòi đúng mã \texttt{911} ở thanh ghi 0x8140 \\
\hline
Loa rè, USB của board rớt 3 lần khi phát to & Bộ khuếch đại lấy 5~V từ chân board, chân này đi từ cáp USB và không tải nổi đỉnh 600~mA & Cấp từ đường 5~V riêng, tụ 470~$\mu$F gần chân nguồn \\
\hline
Đọc dòng quét của panel luôn trả 0xFFFF trong firmware & \texttt{esp\_lcd} tắt driver chân DC giữa hai giao dịch nên lệnh đọc thành dữ liệu & Bật lại ngõ ra DC trước mỗi lần đọc \\
\hline
\end{tabular}
\end{table}

\section{Huấn luyện, lượng tử hóa và triển khai mô hình}
\label{section:4.3}

\subsection{Nguồn của ba mô hình}
\label{section:4.3.1}

Chỉ nhánh phát hiện do đồ án tự huấn luyện. Hai nhánh còn lại nhập trọng số công khai, vì
các bản tự huấn luyện đều thua sau lượng tử hóa hoặc thua trên ảnh OV5640
(Mục~\ref{section:2.3}). Bảng~\ref{tab:nguon_mo_hinh} ghi rõ nguồn và giấy phép. Chuỗi này có
mắt xích chỉ dùng cho nghiên cứu, nên muốn thương mại hóa phải huấn luyện lại nhánh nhận
diện trên dữ liệu khác.

\begin{table}[!htb]
\centering
\small
\caption{Nguồn trọng số và giấy phép của ba mô hình.}
\label{tab:nguon_mo_hinh}
\begin{tabular}{|L{2.5cm}|L{3.8cm}|L{3.6cm}|L{3.3cm}|}
\hline
\textbf{Nhánh} & \textbf{Trọng số} & \textbf{Dữ liệu huấn luyện} & \textbf{Giấy phép} \\
\hline
Phát hiện & Tự huấn luyện YuNet~\cite{yunet2023}, 300 epoch & WIDER FACE~\cite{wider}, 12.880 ảnh, giữ lại 1.262 ảnh để đo & Mã nguồn MIT \\
\hline
Chống giả mạo & Nhập MiniFASNetV1SE của Minivision~\cite{minifas} & Dữ liệu riêng của Minivision & Apache-2.0 \\
\hline
Nhận diện & Nhập MobileFaceNet-ECA của FRBench~\cite{frbench} & MS-Celeb-1M~\cite{ms1m} & Mã nguồn MIT, trọng số chỉ dùng nghiên cứu \\
\hline
\end{tabular}
\end{table}

Dữ liệu đi qua ba tầng: thư mục \texttt{raw/} chỉ đọc, \texttt{interim/} đã chuẩn hóa, và tệp
chia tập ghi kèm seed và sha256. Mỗi lần huấn luyện ghi vào một thư mục mang ngày, mã commit
và băm cấu hình, kèm cấu hình đã phân giải, \texttt{split.lock} và \texttt{env.txt}, nên một
con số trong sổ đo truy được về đúng lần chạy. Script nhập trọng số viết ra thư mục cùng cấu
trúc đó, nên lượng tử hóa không phân biệt mô hình tự huấn luyện hay nhập.

\subsection{Huấn luyện YuNet}
\label{section:4.3.2}

Bản dựng trong đồ án có 75.631 tham số, đầu vào 160$\times$120. Lần huấn luyện chạy 300 epoch
trên RTX 3050 Laptop, batch 64, SGD Nesterov với tốc độ học 0,02, lịch cosine có 3 epoch khởi
động, EMA 0,999, seed 42. Chỗ khó nằm ở công thức lấy mẫu. Ở đầu vào 160$\times$120, ảnh WIDER
bị ép xuống hệ số 0,156, nên trung vị cạnh mặt còn 2,9~px và 85,5\% số mặt nhỏ hơn một ô lưới
bước 8. Tầng bước 32 không nhận mẫu dương nào, tức ba đầu ra chỉ dùng thật một.

Đồ án đổi ba thứ, mọi ngưỡng lấy từ đo: (i)~cắt vùng ngẫu nhiên tỉ lệ 0,3--1,0 rồi phóng về đầu
vào, nên trung vị lên 13,6~px và phủ được dải mặt của máy; (ii)~bỏ mặt dưới 8~px sau khi cắt
chứ không bỏ khỏi tập dữ liệu, để cùng một mặt vẫn được học khi rơi vào vùng cắt gần; (iii)~chia
ba tầng ở các mốc 16 và 48~px, được 61,1\%, 33,6\% và 5,3\% số mặt. Bản FP32 chọn theo AP đạt
0,9414 trên mặt từ 38~px, qua mốc 0,90 của NF3.

\subsection{Nhập trọng số}
\label{section:4.3.3}

Script \texttt{import\_minifasnet.py} dựng lại kiến trúc rồi gấp các khác biệt tiền xử lý vào
trọng số, để code trên chip không phải biết gì về kho gốc. Kho gốc đọc ảnh BGR thang 0--255 nên
script đảo trục kênh vào và nhân 255 vào tích chập đầu, rồi hoán vị đầu ra để lớp sống đứng ở
chỉ số 0. Nhờ vậy mô hình nhận ảnh RGB thang 0--1 đúng như đường tiền xử lý sẵn có.

Script \texttt{import\_frbench.py} nạp trọng số MobileFaceNet-ECA bằng \texttt{load\_state\_dict}
chế độ chặt, giữ đúng tên thuộc tính của kho gốc. Chỗ duy nhất khác kho gốc là khối ECA. Bản gốc
gọi \texttt{expand\_as} trước khi nhân, ONNX ghi bước đó thành phép toán \texttt{Expand}, mà ESP-DL
3.3.11 không có module tương ứng và cũng không có đường dự phòng: nạp bản đầu tiên ngày 23/09/2026
làm máy khởi động lặp 129 lần. Phép nhân vốn tự phát sóng trên kênh, nên bỏ \texttt{expand\_as}
ra cùng hàm, đo trên máy tính lệch 0,0. Từ đó \texttt{espdl\_op\_check.py} đối chiếu phép toán của
mô hình với đúng bản esp-dl trước khi nạp.

\subsection{Lượng tử hóa và tối ưu kernel}
\label{section:4.3.4}

Mỗi nhánh đi theo chuỗi ở Hình~\ref{fig:chuoi_quant}, chia bốn bước. Mỗi bước là một lệnh
\texttt{make} có điều kiện qua, chưa qua thì không sang bước sau.

Bước đúng kiểm chip tính như mô phỏng. \texttt{make quant-timing} dựng một đồ thị thử, rồi
đồ thị đó được nạp lên board bằng \texttt{make ai-unit}. Bước này qua khi \texttt{model->test()} khớp ở cả đồ thị
lẫn các bản cắt giữa đồ thị, và đồ thị khớp mô phỏng từng bit trên khung OV5640 thật. Lệch thì chia
đôi bằng các bản cắt tới phép toán đầu tiên lệch rồi sửa mô phỏng.

Bước nhanh lặp \texttt{make speed}: chạy bản dò trên board, in thời gian từng module, rồi
\texttt{espdl\_profile.py} gắn nhãn các đường chậm kèm bản vá xử lý nó. Bản vá được giữ khi bản dò
vẫn khớp từng bit và thời gian giảm. Bước này dừng khi bảng xếp hạng chỉ còn nhãn không có bản vá.

Bước chất lượng dựng bốn bậc của thang bằng \texttt{quant-ptq}, \texttt{quant-int16},
\texttt{quant-qat}, dựng các hàng có bản vá làm tròn bằng \texttt{quant-variants}, rồi
\texttt{quant-choose} áp luật chọn. Bước chốt chạy lại bản dò trên hàng đã chọn, và
\texttt{quant-deploy} chép tệp vào \texttt{firmware/models/} rồi khóa phiên bản trong
\texttt{models.lock.json}.

Bản dò ghi đè board, nên mỗi phiên mở bằng \texttt{make board-save} và đóng bằng
\texttt{make board-restore}. NVS không bao giờ được lưu hay trả lại, vì một số lần khởi động cũ
làm mã lượt lặp (Mục~\ref{section:3.5.4}).

Bước nhanh sinh ra bảy phép toán tự viết ở Bảng~\ref{tab:phep_toan}. Mỗi phép toán là một cặp: bản
vá Python thay một mẫu trên đồ thị lúc xuất tệp \texttt{.espdl}, module C++ trong
\texttt{ai\_engine} đăng ký tên đó với ESP-DL và chạy nó, phần tính nặng viết bằng lệnh PIE. Hai đầu
khớp nhau qua tên phép toán, tên thuộc tính, cách xếp trọng số và phép làm tròn. Phần mô phỏng phải
chép theo mã máy của kernel ESP-DL đọc bằng \texttt{objdump}, vì esp-dl build với
\texttt{-ffast-math -O3} nên trình biên dịch đã đổi thứ tự các phép số thực. Nâng esp-dl khỏi bản
3.3.11 là phải soát lại toàn bộ các hợp đồng này.

\begin{table}[!htb]
\centering
\small
\caption{Các phép toán đồ án tự viết cho ESP-DL.}
\label{tab:phep_toan}
\begin{tabular}{|L{4.6cm}|L{3.8cm}|L{2.9cm}|L{1.9cm}|}
\hline
\textbf{Phép toán} & \textbf{Đường chậm bị bỏ} & \textbf{Số đo} & \textbf{Nhánh} \\
\hline
\texttt{ConvSlices} & Tích chập 1$\times$1 có trọng số từ 64~KB đọc lại từ PSRAM ở mỗi điểm & Nhận diện 299,1 $\to$ 264,5~ms & Giả mạo, nhận diện \\
\hline
\texttt{ChannelConv} & Năm phép toán của khối ECA & 15 phép còn 0,7~ms & Nhận diện \\
\hline
\texttt{ConvIm2col} & Tích chập đầu 3 kênh đi đường không căn 16 & Tích chập đầu 17,4 $\to$ 8,5~ms & Cả ba \\
\hline
\mbox{\texttt{PointwiseDepthwise}} & Tensor giữa 1$\times$1 và depthwise đi một vòng qua PSRAM & Nhanh hơn 5,0--5,5\% & Phát hiện \\
\hline
\texttt{MulAdd} & Tích của cổng ECA, SE ghi ra PSRAM rồi mới cộng & 12 phép 12,9~ms, trước đó 16,6~ms & Giả mạo, nhận diện \\
\hline
\texttt{SpatialMean} & \texttt{ReduceMean} cộng bằng vòng C tổng quát & 3 phép trước đó 1,9~ms & Giả mạo \\
\hline
\texttt{ConvPRelu} & PReLU chạy riêng sau tích chập & Trượt cửa chất lượng & Không dùng \\
\hline
\end{tabular}
\end{table}

Thang chọn hàng \texttt{mse+clip\_as\_relu} cho phát hiện, \texttt{int16\_top1} cho chống giả mạo và
\texttt{qat} cho nhận diện. Nhánh nhận diện không có hàng nào qua cửa nên luật không chọn được, và đồ
án lấy hàng chất lượng cao nhất. Hàng này đổi số nguyên của mô hình, nên mọi khuôn mặt phải đăng ký
lại sau khi nạp. Số của từng bậc nằm ở Chương~\ref{chapter:Experiment}. Đối chứng TFLite Micro dùng
cùng ba lần huấn luyện, lượng tử hóa ở mức Q1 và chạy bằng một profile riêng.

\subsection{Đưa mô hình lên chip}
\label{section:4.3.5}

\texttt{make pack} gói ba tệp \texttt{.espdl} thành một ảnh có tiêu đề và ghi vào phân vùng
\texttt{models\_0} hoặc \texttt{models\_1}. Lúc khởi động, \texttt{ai\_engine\_init} ánh xạ phân
vùng bằng mmap, không tốn RAM, rồi ESP-DL dựng từng mô hình qua ba bước: đọc đồ thị, chép trọng số
sang PSRAM, lập kế hoạch bộ nhớ tensor. Trọng số được chép vì PSRAM chạy bus 8 dây còn flash chạy 4
dây: đo ngày 23/09/2026, chép sang làm mô hình nhanh hơn từ 3\% tới 25\%. Tensor nào không sống cùng
lúc thì xếp chồng lên cùng một chỗ, nên mỗi mô hình chỉ cần một khối. Bảng~\ref{tab:bo_nho_mo_hinh}
là kết quả đo trên board.

\begin{table}[!htb]
\centering
\small
\caption{Bộ nhớ và thời gian dựng của ba mô hình trên board.}
\label{tab:bo_nho_mo_hinh}
\begin{tabular}{|L{3.0cm}|c|c|c|c|}
\hline
\textbf{Nhánh} & \textbf{Tệp (KB)} & \textbf{Trọng số (KB)} & \textbf{Tensor (KB)} & \textbf{Dựng (ms)} \\
\hline
Phát hiện & 192,9 & 78,4 & 187,5 & 62 \\
\hline
Chống giả mạo & 577,1 & 464,2 & 568,8 & 123 \\
\hline
Nhận diện & 1.345,8 & 1.194,6 & 820,8 & 192 \\
\hline
\end{tabular}
\end{table}

Cả ba lấy 3.574~KB PSRAM và 0~KB RAM nội. Phần chênh so với tổng của bảng là hàng nghìn đối tượng
nhỏ của ESP-DL, khoảng 73--89~KB mỗi mô hình. Một guard đẩy chúng sang PSRAM trong lúc dựng, vì để
mặc định thì ba mô hình lấy 94~KB RAM nội, phần bộ nhớ mà Wi-Fi, TLS và DMA đang cần.

\section{Xây dựng firmware}
\label{section:4.4}

\subsection{Profile build}
\label{section:4.4.1}

Firmware dựng theo bốn profile ở Bảng~\ref{tab:profile}, mỗi profile trong thư mục riêng. Số đo chỉ
lấy ở \texttt{bench}. Đo ngày 14/10/2026, cùng mã nguồn dựng bằng \texttt{dev} ghi một khung màn hình
mất 46,0--55,0~ms và xé hình, còn \texttt{bench} mất 36,9--39,5~ms và không xé.

\begin{table}[!htb]
\centering
\small
\caption{Các profile build của firmware.}
\label{tab:profile}
\begin{tabular}{|L{1.6cm}|L{1.5cm}|L{10.4cm}|}
\hline
\textbf{Profile} & \textbf{Tối ưu} & \textbf{Đặc điểm} \\
\hline
\texttt{dev} & \texttt{-Og} & Gỡ lỗi: assert, kiểm tràn stack, kiểm heap, gdbstub, console \\
\hline
\texttt{bench} & \texttt{-O2} & Tốc độ như bản xuất xưởng, còn trace FreeRTOS và console \\
\hline
\texttt{prod} & \texttt{-O2} & Log mức cảnh báo, tắt console, bắt buộc TLS, bắt buộc có token lô \\
\hline
\texttt{fleet} & \texttt{-O2} & Bản phát hành OTA: \texttt{prod} cộng cấu hình khởi động lại khi panic \\
\hline
\end{tabular}
\end{table}

Secure Boot và mã hóa flash chưa bật ở profile nào, vì đốt eFuse là thao tác một chiều. Người có
máy trong tay vì thế đọc được JWT từ flash và gửi được lượt chấm cho người đang ở máy đó. Máy chủ chặn
một phần bằng cặp máy, người và luật giữ lượt ở Mục~\ref{section:3.5.5}, phần còn lại là việc chưa
làm.

\subsection{Driver ngoại vi}
\label{section:4.4.2}

Camera chạy một cấu hình duy nhất cho xem trước, suy luận và thu dữ liệu: khung 480$\times$320
RGB565, XCLK 27~MHz, 14,19~fps. Đồ án tắt tự động phơi sáng của OV5640 vì nó đo sáng cả cảnh, trần
nhà sáng dìm mặt vào bóng: đo lúc đó, điểm sống của cùng một mặt đi từ 0,9999 xuống 0,683. Hàm
\texttt{drv\_camera\_expose} tự đo một phần tư giữa khung, mỗi lần chỉnh một phần tư quãng còn thiếu
rồi nghỉ hai khung cho cảm biến bắt kịp. Thời gian phơi lấy theo bội 10~ms để không nhấp nháy dưới
đèn 50~Hz, phần lẻ bù bằng gain 1--4 lần.

Ghi một khung màn hình mất 31,7~ms ở SPI 80~MHz, lâu hơn chu kỳ quét 57~Hz của panel, nên tia quét
vượt con trỏ ghi và ảnh động bị khấc ngang. \texttt{drv\_lcd} sửa bằng hai việc. Thanh ghi
\texttt{FRMCTR1} hạ tần số quét xuống 23,26~Hz, chu kỳ 42,98~ms, dài hơn thời gian ghi. Trước mỗi lượt
ghi, driver đọc dòng quét qua chân SDO và chỉ bắt đầu khi bộ đếm ở đoạn 160--200, nên con trỏ ghi có
ít nhất 7,2~ms chạy đà. Tốc độ khung của ảnh không đổi, vì camera mới là thứ quyết định nhịp đó. Lần
đọc đầu ngay sau một lượt ghi luôn ra giá trị ngoài dải, nên driver đọc lại tối đa 4 lần. Bỏ cuộc ngay
lần đầu làm 616/1.260 khung ghi lệch nhịp mà không báo gì.

Các driver còn lại mỏng hơn. \texttt{drv\_touch} chạy trình tự chọn địa chỉ GT911 rồi đọc điểm chạm theo
ngắt. \texttt{drv\_tof} đặt cực tính ngắt và trả khoảng cách, task \texttt{tof} đổi nó thành hai cạnh có
người và hết người. \texttt{drv\_audio} nạp tiếng bíp tổng hợp từ phân vùng \texttt{assets} một lần lúc
khởi động, giữ PCM trong PSRAM, và kéo chân SD của bộ khuếch đại xuống khi không phát. Mọi lệnh tới
PCF8574 và các thiết bị khác trên bus đi qua khóa \texttt{m\_i2c}.

\subsection{Đường đi của một khung}
\label{section:4.4.3}

Hình~\ref{fig:duong_khung} là đường đi của một khung. Hàng đợi chỉ chở con trỏ 4 byte, nên 300~KB
điểm ảnh nằm yên ở chỗ DMA đã ghi và ba mô hình cùng đọc tại đó. Mỗi lúc chỉ một task giữ khung nên
không cần khóa. Cái giá là AI giữ khung suốt lượt chạy, tới khoảng 350~ms, nên camera cần 4 bộ đệm: một
cho AI, một nằm chờ trong hàng đợi, một cho \texttt{cam} đang vẽ, một cho cảm biến đang ghi. Với 3 bộ
đệm, xem trước tụt từ 14,18~fps xuống 8,1~fps.

\begin{figure}[!htb]
\centering
\begin{tikzpicture}[
  font=\footnotesize,
  node distance=9mm and 14mm,
  box/.style={draw, align=center, text width=24mm, minimum height=11mm, inner sep=2pt},
  lab/.style={font=\scriptsize, align=center},
  arr/.style={-{Latex[length=2mm]}}]
\node[box] (s) {OV5640};
\node[box, right=of s] (b) {Bốn bộ đệm\\ở PSRAM};
\node[box, right=of b] (c) {Task \texttt{cam},\\lõi 0};
\node[box, right=of c] (p) {Panel ST7796S};
\node[box, below=of c] (a) {Task \texttt{ai},\\lõi 1};
\node[box, right=of a] (t) {Tensor vào của\\từng mô hình};
\node[box, left=of a] (r) {Task \texttt{attend}};
\draw[arr] (s) -- node[lab, above] {DMA} (b);
\draw[arr] (b) -- (c);
\draw[arr] (c) -- node[lab, above] {SPI} (p);
\draw[arr] (c) -- node[lab, right] {\texttt{q\_frame\_ai}, con trỏ} (a);
\draw[arr] (b.south east) -- node[lab, sloped, below] {đọc tại chỗ} (a.north west);
\draw[arr] (a) -- node[lab, below] {cắt,\\thu nhỏ} (t);
\draw[arr] (a) -- node[lab, below] {\texttt{q\_result}} (r);
\end{tikzpicture}
\caption{Đường đi của một khung trong firmware.}
\label{fig:duong_khung}
\end{figure}

Thứ duy nhất bị chép là bản đã cắt hoặc thu nhỏ, ghi thẳng vào tensor vào của từng mô hình: 57,6~KB,
19,2~KB và 37,6~KB. \texttt{letterbox\_frame} thu cả khung về 160$\times$120, \texttt{crop\_face} cắt
vùng 2,7 lần hộp mặt từ khung gốc. \texttt{align\_face} tìm phép xoay, phóng đều và dời từ năm điểm mốc
bằng bình phương tối thiểu dạng đóng, rồi đi từ từng điểm của ảnh đích 112$\times$112 ngược về khung gốc
và nội suy song tuyến, nên ảnh đích không có lỗ. Vector ra được chuẩn hóa về độ dài 1 rồi lượng tử lại
về int8 với thang riêng từng vector. Phép so trong \texttt{svc\_facedb} dùng nhân vô hướng int8 viết bằng
assembly PIE và chia cho bình phương độ dài đã tính sẵn của mẫu.

\texttt{VisionPipeline} nhận bốn giao diện cho phát hiện, chống giả mạo, nhận diện và so khớp. Nhờ
vậy ứng dụng kiểm thử cắm bốn bản giả đếm số lần gọi, chạy trên board mà không cần
phân vùng mô hình, và kiểm đúng thứ tự ba mô hình cùng mọi lối thoát sớm.

\subsection{Máy trạng thái và đường lượt chấm}
\label{section:4.4.4}

Máy trạng thái ở Hình~\ref{fig:may_trang_thai} là một mảng \texttt{constexpr} 24 dòng, hàm
\texttt{next} tra mảng và trả trạng thái kế tiếp cùng hành động. Hai đường tự lặp của \texttt{Granted}
và \texttt{Denied} không đổi trạng thái, nên \texttt{main} không thể đoán từ cạnh trạng thái. Hàm
\texttt{svc\_attendance\_on\_vision} vì thế trả thẳng một trong bốn lời đáp \texttt{GRANTED},
\texttt{ALREADY}, \texttt{REFUSED}, \texttt{NOTHING}, và màn hình vẽ theo đó.

Mỗi lượt chấm ghi bằng \texttt{lfs\_file\_write} rồi \texttt{lfs\_file\_sync} ngay, tốn 10--20~ms, xong
mới đưa lời nhắc vào \texttt{q\_uplink}. Task \texttt{sync} đọc bản ghi ở con trỏ, dựng JSON bằng kiểu sinh
từ hợp đồng, gửi QoS 1, chờ PUBACK rồi mới ghi \texttt{cursor.bin} theo đường hai pha. Bản ghi làm ra lúc
máy chưa có giờ được cộng mốc khởi động trước khi gửi, và con trỏ dừng lại chờ giờ tối đa mười phút, vì
broker thường lên trước NTP.

RAM nội là giới hạn thật của phần mạng. Đo trên \texttt{dev}, mảnh liền lớn nhất chỉ còn 5.620~B trong khi
esp-mqtt xin 6.144~B cho ngăn xếp, nên ngăn xếp task MQTT nằm ở PSRAM và mbedTLS cũng cấp phát ở đó. Điều
kiện của cách này là ngăn xếp ở PSRAM không được chạm khi cache flash tắt, và lượt ghi flash đã đỗ lõi kia
và tắt ngắt nên không task nào chạy lúc đó.

\subsection{OTA}
\label{section:4.4.5}

\texttt{net\_ota} so link với gốc API đã biên dịch sẵn trước khi mở kết nối, tải vào khe không chạy, kiểm
sha256 rồi đổi khe. Khởi động được chưa đủ để coi là thành công, vì một gói mô hình hỏng vẫn boot được rồi
không nhận ra ai. Bản mới ở trạng thái chạy thử tới khi chạy đủ 30~s và \texttt{AI\_READY} đã bật, lúc đó
hàm \texttt{settle\_this\_build} mới đánh dấu khe hợp lệ cho bootloader. Gói mô hình đi
đường tương tự, với dấu chạy thử nằm trong NVS: ghi dấu trước, đổi khe sau, và gói không nạp được thì máy
lùi khe rồi khởi động lại. Nhịp sống mang cờ \texttt{onTrial}, nên dashboard phân biệt được bản đã lên với bản
đang chạy thử.

\subsection{Giao diện máy chấm công}
\label{section:4.4.6}

Mỗi màn hình là một lớp con của \texttt{Screen}, \texttt{ScreenManager} giữ chín đối tượng dựng sẵn lúc khởi
động. \texttt{ui} vẽ vào \texttt{Canvas}, một bản đồ 1 byte mỗi điểm mang chỉ số màu và độ phủ, rồi \texttt{cam}
trộn nó vào từng dải 20 dòng ngay lúc đẩy ra panel, nên lớp phủ không tốn thêm lượt SPI hay lượt quét PSRAM
nào. Hai khoang của lớp phủ cần một luật giữ chỗ: khi \texttt{ui} được phép xóa khoang \texttt{cam} đang đọc,
đo 19/09/2026 có 156/660 khung hỏng lúc chạm liên tục, thêm luật giữ chỗ thì 0/840 khung.

Chữ lấy từ bảng glyph 4~bit sinh sẵn từ phông MiSans, khoảng 110 chuỗi khai ở \texttt{StrId} với hai bản dịch.
Hình~\ref{fig:man_may} là màn quét và màn lấy mẫu. Màn lấy mẫu vẽ ba vạch sáng dần, đổi câu nhắc theo mẫu đang
chờ, và giữ tên người đang được chụp ở thanh trên để mặt không rơi vào tên người khác. Màn Cài đặt
chia thẻ nhóm gồm ngôn ngữ, Wi-Fi, độ sáng, âm lượng và trang thiết bị của tôi, giá trị ghi vào NVS.

\begin{figure}[!htb]
\centering
\anh{anh_man_quet.jpg}{0.4\textwidth}{7cm}
\hspace{4mm}
\anh{anh_man_mau.jpg}{0.4\textwidth}{7cm}
\caption{Màn quét và màn lấy mẫu trên máy chấm công.}
\label{fig:man_may}
\end{figure}

\subsection{Kiểm thử trên board}
\label{section:4.4.7}

Firmware có 29 ứng dụng kiểm thử, 25 ứng dụng nằm trong từng thành phần và 4 ứng dụng chung
\texttt{bench\_ai}, \texttt{bench\_mem}, \texttt{parity}, \texttt{soak}, tổng 266 ca Unity. \texttt{make fw-app-test} nạp
ứng dụng và chạy kịch bản pytest trên board. Bộ vector vàng ở \texttt{contracts/golden/} giữ hậu xử lý C++
khớp bản Python. Một dòng tổng kết màu xanh chưa đủ để tin: bộ kiểm thử lưu trữ từng in 11/11 ca đạt trong khi
phía trên có hai lần panic, nên kịch bản đếm cả số lần khởi động.

\section{Xây dựng backend và các phân hệ nhân sự}
\label{section:4.5}

\subsection{Cấu trúc}
\label{section:4.5.1}

Backend là một ứng dụng NestJS~\cite{nestjs} gồm 27 module, nối PostgreSQL qua Prisma~\cite{prisma}. Phần dùng
chung nằm ở \texttt{src/common/}: guard xác thực và vai, \texttt{ScopeService}, interceptor ghi nhật ký và phát tin
thay đổi, bộ ghi CSV. Mọi biến môi trường đọc ở một tệp \texttt{env.schema.ts} và kiểm bằng zod lúc khởi động,
thiếu là dừng ngay. Mỗi route khai vai bằng \texttt{@Roles} hoặc \texttt{@CurrentViewer}, và
\texttt{check\_routes.py} chặn route thiếu khai báo. \texttt{check\_error\_codes.py} kiểm mọi mã lỗi có câu ở
cả hai ngôn ngữ. Hình~\ref{fig:luong_cong} là đường dữ liệu từ lượt chấm tới phiếu lương, các mục dưới đi theo
hình này.

\begin{figure}[!htb]
\centering
\begin{tikzpicture}[
  font=\footnotesize,
  node distance=6mm and 7mm,
  box/.style={draw, align=center, text width=30mm, minimum height=11mm, inner sep=2pt},
  arr/.style={-{Latex[length=2mm]}}]
\node[box] (r) {AttendanceRecord\\lượt chấm thô};
\node[box, right=of r] (d) {AttendanceDay\\ngày công};
\node[box, right=of d] (p) {PayrollRun\\lượt tính};
\node[box, right=of p] (s) {Payslip\\phiếu lương};
\node[box, below=of d] (q) {Request\\đơn nghỉ, giải trình};
\node[box, below=of p] (c) {PayrollPolicy,\\mức lương};
\draw[arr] (r) -- (d);
\draw[arr] (d) -- (p);
\draw[arr] (p) -- (s);
\draw[arr] (q) -- (d);
\draw[arr] (c) -- (p);
\end{tikzpicture}
\caption{Đường dữ liệu từ lượt chấm tới phiếu lương.}
\label{fig:luong_cong}
\end{figure}

\subsection{Nhận lượt chấm qua MQTT}
\label{section:4.5.2}

\texttt{MqttService} nối EMQX bằng client ID cố định, phiên không sạch, và tự giữ việc trả ack: thư viện mqtt.js
trả ack ngay trong vòng đọc, còn đồ án hoãn tới khi trình xử lý xong. Tin của mỗi máy đi theo một làn riêng nên
đúng thứ tự tới, các máy không chờ nhau. Tin hỏng được thử lại giãn từ 1~s tới 30~s, sau năm lần thì ghi log rồi
ack, để một tin độc không chặn cả hàng.

\texttt{AttendanceListener} so \texttt{deviceId} trong bản tin với topic, vì topic mới là thứ broker đã xác thực.
\texttt{AttendanceService.record} rồi kiểm lần lượt: máy còn trong đội, lượt chưa có, người đang ở máy đó, giờ có
nghi vấn không, có phải giữ lượt không. Hai lần giao cùng lúc có thể cùng qua phép đọc, và ràng buộc unique ở cơ
sở dữ liệu là chỗ phân xử cuối. Lượt của một ngày đã qua xếp ngày đó của người đó vào hàng dựng lại. Ghi xong mới
phát tin cho dashboard qua Socket.IO~\cite{socketio}, vì dashboard trả lời bằng cách hỏi lại danh sách.

\subsection{Ngày công, ca làm và giải trình công}
\label{section:4.5.3}

Việc dựng ngày chạy lúc 00:30 theo múi giờ công ty cho ngày hôm qua. Hàm \texttt{dayOf} lấy lượt đầu và lượt cuối
của ngày, tính phút làm, phút muộn sau ân hạn, phút về sớm và phút tăng ca so với ca của người đó. Ngày không có
lượt thì tra lịch: ngày lễ, cuối tuần, phép đã duyệt, công tác đã duyệt, còn lại là vắng. Phút tăng ca chỉ được đếm,
chỉ phần có đơn duyệt mới thành tiền. Nghỉ nửa ngày dời một đầu của ca về điểm giữa ca, ví dụ ca 08:00--17:30 có
điểm giữa 12:45. Ngày nhân sự đã sửa tay không bị dựng đè, và dòng ngày giữ người sửa, lý do, phút máy đo.

Ngày lệch sinh một việc cho chính người đó. Từ đó họ gửi đơn giải trình công, là một \texttt{Request} loại
\texttt{ATTENDANCE\_FIX} điền sẵn ngày ấy, hoặc bấm không cần giải trình. Đơn được duyệt thì \texttt{applyFix} sửa
ngày công trong cùng giao dịch với lượt duyệt. Ca qua nửa đêm chưa có thiết kế: ca có giờ ra không đứng sau giờ vào
bị từ chối với mã \texttt{SHIFT\_CROSSES\_MIDNIGHT}, vì mọi phép tính trên đặt cả ca vào một ngày lịch.

\subsection{Nghỉ phép}
\label{section:4.5.4}

Số dư giữ chỗ bằng một câu \texttt{UPDATE} có điều kiện, kiểm và giữ trong cùng một lệnh, vì hai đơn đọc cùng một số
dư sẽ cùng qua phép kiểm tách rời. Câu lệnh không chạm dòng nào thì trả \texttt{LEAVE\_BALANCE\_SHORT}. Đơn chồng ngày
do ràng buộc loại trừ của PostgreSQL chặn. Lượt duyệt cũng là ghi có điều kiện \texttt{state = PENDING}, nên hai người
duyệt cùng lúc thì người sau nhận \texttt{REQUEST\_ALREADY\_DECIDED}. Đơn mang \texttt{clientKey} do máy gửi sinh ra,
nên điện thoại gửi lại lúc mạng yếu chỉ nhận về đơn đã có, không thành hai đơn.

\subsection{Tính lương}
\label{section:4.5.5}

Phép tính nằm trong một hàm thuần \texttt{calculate}, không đụng cơ sở dữ liệu nên kiểm thử được từng ca. Tiền là số
nguyên lớn tính bằng đồng, tỉ lệ là phần vạn, ngày là phần trăm ngày. Thứ tự tính gồm: (i)~lương theo ngày công;
(ii)~phụ cấp, tăng ca theo từng hệ số, phụ trội làm đêm; (iii)~bảo hiểm trên lương đóng bảo hiểm, có trần, và miễn
đóng khi nghỉ không lương từ 14 ngày; (iv)~thu nhập tính thuế sau phần miễn, bảo hiểm và giảm trừ gia cảnh; (v)~thuế
lũy tiến; (vi)~các khoản trừ khác như tạm ứng, rồi thực lĩnh. Phần âm không thành dòng chuyển tiền âm mà ghi thành
dòng \texttt{CARRY\_OUT} và sang kỳ sau.

Mọi con số do luật đặt nằm ở \texttt{PayrollPolicy} và \texttt{TaxBracket} theo ngày hiệu lực, và phép tính hỏi chính
sách tại ngày cuối kỳ, nên mức tham chiếu đổi giữa năm 2026 là hai dòng dữ liệu. Bộ kiểm thử kiểm biểu thuế bằng hai
đường độc lập: liên tục tại mọi ranh giới bậc, và khớp công thức rút gọn của từng bậc.

Bấm chạy thì \texttt{execute} kiểm trước những gì chắc chắn làm lượt hỏng, rồi xếp việc vào hàng đợi. \texttt{runNow}
chia 500 người một lô, mỗi lô một giao dịch, nên không có phiếu nào thiếu dòng. Khóa đặt trên dòng pháp nhân, dòng kỳ
và dòng lượt, nên hai tài khoản bấm cùng lúc gặp đúng một luật. Một điểm còn lệch: hàm này chỉ miễn phần chênh của tiền
làm thêm giờ theo Thông tư 111/2013, trong khi văn bản hướng dẫn Luật Thuế thu nhập cá nhân
2025~\cite{luattncn2025} được cho là miễn toàn bộ. Đồ án chưa đọc được nguyên văn nên chưa sửa code.

\subsection{Quản lý máy chấm công và hàng đợi}
\label{section:4.5.6}

Module \texttt{devices} giữ vòng đời máy, nhận nhịp sống, ghi độ lệch đồng hồ và gửi lệnh qua \texttt{down/cmd} với thời
hạn. Module \texttt{enrollment} giữ cặp máy, người và đẩy danh sách xuống theo số phiên bản. Module \texttt{models} nhận
bản phát hành qua \texttt{POST /releases}, băm sha256 trong lúc ghi xuống volume, và tìm đường ngắn nhất từ thứ máy đang
chạy tới bản mới nhất của cả firmware lẫn mô hình, theo các lời khai \texttt{minFwVersion} và \texttt{requiresModels}.

Việc nền đi qua sáu hàng đợi BullMQ~\cite{bullmq}: \texttt{report}, \texttt{notify}, \texttt{payroll},
\texttt{timesheet}, \texttt{people}, \texttt{leave}, gồm 31 loại việc. Mọi lịch lặp mang múi giờ công ty, và dựng lại một
ngày không đổi gì thì không ghi dòng nào. Thông báo tách tin của một người khỏi việc của cả nhóm: ba người bàn nhân sự
cùng nhận một đơn, một người duyệt xong thì việc đóng cho cả ba.

\section{Xây dựng dashboard, cổng nhân viên và ứng dụng Android}
\label{section:4.6}

\subsection{Dashboard}
\label{section:4.6.1}

Frontend là một ứng dụng Next.js~\cite{nextjs} dùng App Router, 46 trang nằm dưới tiền tố ngôn ngữ, giao diện dựng bằng
Kumo~\cite{kumo}. Dữ liệu đi qua TanStack Query~\cite{tanstack}, bộ lọc giữ trong URL, và \texttt{lib/ws.ts} ánh xạ mỗi
kênh tin sang các truy vấn cần làm mới, nên trang đang mở tự cập nhật. Chuỗi hiển thị dùng next-intl với 2.993 khóa mỗi
ngôn ngữ. Tệp \texttt{vi.json} là nguồn kiểu nên \texttt{en.json} thiếu khóa là lỗi biên dịch, và \texttt{lib/fault.ts} đổi mã
lỗi của API thành câu.

Các ảnh dưới chụp từ môi trường thử với một công ty giả 5.006 nhân viên, tên người là tên sinh ngẫu nhiên.
Hình~\ref{fig:web_duyet} là hộp chờ duyệt: mỗi dòng đủ ngữ cảnh để quyết, gồm số dư sau khi duyệt và thời gian đã chờ, số
đếm trên tab đến từ tổng của máy chủ. Hình~\ref{fig:web_luong} là một kỳ lương đã trả với 4.767 phiếu.
Hình~\ref{fig:web_may} là danh sách máy chấm công với trạng thái, phiên bản và máy chờ duyệt.

\begin{figure}[!htb]
\centering
\includegraphics[width=0.88\textwidth]{web_cho_duyet.png}
\caption{Hộp chờ duyệt của nhân sự, dữ liệu giả.}
\label{fig:web_duyet}
\end{figure}

\begin{figure}[!htb]
\centering
\includegraphics[width=0.88\textwidth]{web_ky_luong.png}
\caption{Trang một kỳ lương đã trả, dữ liệu giả.}
\label{fig:web_luong}
\end{figure}

\begin{figure}[!htb]
\centering
\includegraphics[width=0.88\textwidth]{web_may_cham_cong.png}
\caption{Danh sách máy chấm công, dữ liệu giả.}
\label{fig:web_may}
\end{figure}

\subsection{Cổng nhân viên}
\label{section:4.6.2}

Cổng nhân viên là chính ứng dụng web ở bề ngang hẹp, Hình~\ref{fig:dien_thoai}. Bảng đổi thành thẻ từ cùng một bộ cột,
mỗi cột mang mức ưu tiên để bảng tự bỏ cột theo bề rộng thật. Ứng dụng cài được như PWA với \texttt{manifest.webmanifest} và
\texttt{sw.js}. Service worker giữ dữ liệu đã tải trong ngăn của từng tài khoản và xóa khi đăng xuất. Đơn gửi lúc mất mạng
nằm trong IndexedDB ở \texttt{lib/outbox.ts} và được thử lại ở ba lúc: khi trình duyệt báo có mạng, khi mở lại ứng dụng, và
một nhịp chậm trong lúc mở. Đồ án không dùng Background Sync vì Safari trên iOS không có. Thông báo đẩy đi qua Web
Push~\cite{webpush} với khóa VAPID, nửa riêng không rời VPS.

\begin{figure}[!htb]
\centering
\includegraphics[width=0.29\textwidth]{dt_trang_chu.png}
\hspace{2mm}
\includegraphics[width=0.29\textwidth]{dt_cham_cong.png}
\hspace{2mm}
\includegraphics[width=0.29\textwidth]{dt_phieu_luong.png}
\caption{Trang của quản lý, công của tôi và phiếu lương trên điện thoại, dữ liệu giả.}
\label{fig:dien_thoai}
\end{figure}

\subsection{Ứng dụng Android}
\label{section:4.6.3}

Trên Android, Chrome không cho PWA đã cài tô hai thanh hệ thống theo màu ứng dụng: đo ngày 25/09/2026 trên điện thoại thật,
hai dải đó lóe khác màu mỗi lần mở. Đồ án vì thế bọc đúng trang web trong một vỏ Trusted Web Activity~\cite{twa}, gói
\texttt{vn.io.cckiosk.nhanluc}, dựng bằng Bubblewrap từ \texttt{twa-manifest.json}. Vỏ không mang code nghiệp vụ nào, nên mọi
thay đổi của web tới ứng dụng như tới trình duyệt. Tệp \texttt{assetlinks.json} trên miền web mang dấu SHA-256 của khóa ký,
thiếu nó thì Chrome hiện thanh địa chỉ. Tệp APK phát qua GitHub Releases, còn iPhone dùng PWA vì iOS không cài APK.

\section{Triển khai Docker và CI/CD}
\label{section:4.7}

\subsection{Dịch vụ trên VPS}
\label{section:4.7.1}

Máy chủ chạy bằng Docker Compose~\cite{compose} gồm Traefik~\cite{traefik}, API, PostgreSQL, Redis, EMQX và một container sao
lưu. Frontend không nằm ở đây mà build trên Vercel mỗi lần push, dưới tên miền \texttt{app.cckiosk.io.vn}, cùng site với
\texttt{api.cckiosk.io.vn}, để cookie làm mới là cookie bên thứ nhất và Safari không đăng xuất người dùng iPhone.

EMQX mặc định mở, nên \texttt{emqx.conf} phải đóng lại ba thứ: xác thực hai tầng, tầng đầu cho tài khoản dịch vụ, tầng sau hỏi
\texttt{POST /mqtt/auth}~\cite{emqxhttp}; ACL theo tên đăng nhập; chứng thư thật do CA riêng ký. Retained mặc định nằm ở RAM, đo
được 0 bản ghi còn lại sau khi khởi động lại broker, nên phải đặt ghi xuống đĩa, nếu không máy đang tắt sẽ trông như chưa từng
tồn tại. Đường \texttt{/mqtt/auth} không ra internet: router của Traefik loại tiền tố đó không phân biệt hoa thường, vì Express
định tuyến không phân biệt mà \texttt{PathPrefix} thì có.

\subsection{Tích hợp và triển khai liên tục}
\label{section:4.7.2}

Bảng~\ref{tab:ci} liệt kê sáu workflow của GitHub Actions~\cite{ghactions}. Workflow \texttt{contracts} chạy mọi lần push và giữ
các khối không trôi khỏi nhau, kể cả sơ đồ nguyên lý so với bảng chân.

\begin{table}[!htb]
\centering
\small
\caption{Các workflow của GitHub Actions.}
\label{tab:ci}
\begin{tabular}{|L{2.4cm}|L{11.2cm}|}
\hline
\textbf{Workflow} & \textbf{Việc} \\
\hline
\texttt{contracts} & Kiểm lược đồ, sinh lại mã và chặn lệch, luật comment, tài liệu, tầng firmware, sơ đồ, mạch in, mã lỗi \\
\hline
\texttt{ml} & Lint và pytest trên CPU \\
\hline
\texttt{firmware} & Build và báo cỡ ảnh; job \texttt{release} phát hành bản \texttt{fleet} khi số phiên bản mới \\
\hline
\texttt{backend} & PostgreSQL, Redis, EMQX thật, migrate, seed, kiểm thử e2e, build \\
\hline
\texttt{frontend} & Sinh kiểu route, \texttt{tsc}, \texttt{next build} \\
\hline
\texttt{deploy} & Gọi lại \texttt{backend}, dựng ảnh container, chạy \texttt{deploy.sh} trên VPS \\
\hline
\end{tabular}
\end{table}

Hình~\ref{fig:cd} là đường triển khai. Push lên \texttt{main} chạm backend, hợp đồng hay cấu hình triển khai là lên VPS, nhưng chỉ
khi kiểm thử qua. Ảnh container dựng trên GitHub vì đo ngày 24/09/2026 trên chính VPS, build mất 3,5 phút với đỉnh 1.250~MB RAM,
trong khi cả stack lúc chạy chỉ khoảng 390~MB. Khóa SSH của GitHub bị ép chạy đúng \texttt{deploy.sh}, nên lộ khóa chỉ deploy lại
được một bản đã có. Token kéo ảnh đi qua stdin và hết hạn cùng job.

\begin{figure}[!htb]
\centering
\begin{tikzpicture}[
  font=\footnotesize,
  node distance=6mm and 5mm,
  box/.style={draw, align=center, text width=28mm, minimum height=11mm, inner sep=2pt},
  arr/.style={-{Latex[length=2mm]}},
  lab/.style={font=\scriptsize, align=center}]
\node[box] (p) {Push lên\\\texttt{main}};
\node[box, right=of p] (t) {Job \texttt{test}};
\node[box, right=of t] (i) {Job \texttt{image}};
\node[box, right=of i] (s) {Job \texttt{ship}:\\\texttt{deploy.sh}};
\node[box, below=of s] (h) {Chờ \texttt{/health}\\tối đa 120~s};
\node[box, left=of h] (d) {Ghi \texttt{.deployed},\\dọn ảnh cũ};
\node[box, below=of h] (r) {Lùi về sha\\ở \texttt{.deployed}};
\draw[arr] (p) -- (t);
\draw[arr] (t) -- (i);
\draw[arr] (i) -- (s);
\draw[arr] (s) -- (h);
\draw[arr] (h) -- node[lab, above] {200} (d);
\draw[arr] (h) -- node[lab, right] {không đạt} (r);
\end{tikzpicture}
\caption{Đường triển khai backend lên VPS.}
\label{fig:cd}
\end{figure}

\texttt{deploy.sh} đưa cấu hình về đúng commit, dựng lại image sao lưu khi tệp của nó đổi, kéo ảnh API, rồi dựng lại bằng
\texttt{-{}-force-recreate} đúng dịch vụ có tệp cấu hình đổi. Bước cuối này cần thiết vì \texttt{git checkout} thay tệp bằng inode mới
trong khi bind mount một tệp lẻ vẫn cầm inode cũ, nên dịch vụ chạy tiếp cấu hình cũ mà không báo gì. Container API chạy
\texttt{prisma migrate deploy} trước khi nghe, và lùi bản an toàn nhờ luật nở rồi co: bản cũ vẫn chạy được trên lược đồ mới.

Firmware và mô hình đi đường khác. \texttt{make release} tăng số phiên bản, dựng ảnh như CI, commit và theo lượt CI tới khi máy chủ có
bản mới. \texttt{make models-release} chạy trên máy huấn luyện, vì trọng số không vào kho mã.

\subsection{Sao lưu}
\label{section:4.7.3}

Mỗi đêm container sao lưu viết một bản logic, một bản logic riêng cho mẫu khuôn mặt và một bản gốc vật lý, còn PostgreSQL đẩy WAL liên
tục với \texttt{archive\_timeout} 5 phút. Hai loại cần nhau vì WAL chỉ phát lại được lên bản gốc vật lý. Tệp được nén rồi mã hóa bằng
khóa công khai của \texttt{age} trước khi chạm đĩa, nên ai chiếm được máy chủ cũng không đọc lại được lịch sử. Mỗi 5 phút tệp mới được
chép sang một bucket ngoài máy, có khóa 5 ngày chống xóa. Script \texttt{restore-drill.sh} dựng lại cả hai loại vào một PostgreSQL vứt đi
và bấm giờ, cứ 15 phút API hỏi trạng thái đẩy WAL và gửi mail cho quản trị khi có chuỗi cũ quá 26 giờ. Trong 7 ngày gần nhất hệ thống mất
tối đa 5 phút dữ liệu, cũ hơn thì mất tối đa một ngày.

Hệ thống đã hiện thực đủ các khối phần mềm ở Chương~\ref{chapter:Design}: firmware 0.9.13 với 29 ứng dụng kiểm thử, ba mô hình trên ESP-DL
với bảy phép toán tự viết, backend 27 module chạy trên VPS từ tháng 9/2026, 46 trang web và ứng dụng Android. Phần cứng còn thiếu mạch đế in,
vỏ và thanh chắn. Chương~\ref{chapter:Experiment} đo các phần trên và đối chiếu với yêu cầu đã đặt.

\end{document}
